\section{Trivalent compression and the exponent}\label{sec:shaving}

Each application of Lemma~\ref{lem:macro} reduces the size per copy and
raises the maximum arity. We stop while the saving in size still exceeds
the cost of evaluating the larger tables. Two operations then recover a
trivalent network over the original field. Batched evaluation uses
Lemma~\ref{lem:batch-table}, and a single coordinate projection removes
the field extension used in the local identities.

\begin{lemma}[Descent of a trivalent network]\label{lem:table-descent}
Let $K/k$ have finite degree $D$. If a trivalent $K$-network of size
$s$ represents a tensor with entries in $k$, the same tensor has a
trivalent $k$-network of size $O(sD^2)$ with the same exterior positions.
\end{lemma}

\begin{proof}
Choose a $k$-basis of $K$ whose first element is $1$. For
$\lambda\in K$, let $M_\lambda$ represent multiplication by $\lambda$
in that basis. Order the stations along a chain and give each of them
two extra ports carrying the incoming and outgoing basis indices.
If an assignment to its original ports has weight $\lambda$, the
enlarged table returns the indicated entry of $M_\lambda$. Pin both
ends of the chain to the first index. For any fixed assignment to the
original wires, contracting along the chain now gives the first
coordinate of the product of the weights. This coordinate map is
$k$-linear. Summing over the original internal bits therefore takes
the first coordinate of the old tensor entry, which belongs to $k$.

Encode each basis index in $\lceil\log_2 D\rceil$ bits, giving weight
zero to unused codes. The enlarged table has at most
$r=3+2\lceil\log_2 D\rceil$ bits. Every table on $r$ bits has a
trivalent realization of size $O(2^r)$. In the binary prefix tree
compute, at each node, the indicator that the input has that prefix.
At depth $j$, copying the next input bit to the $2^j$ nodes and testing
the two children of each node costs $O(2^j)$ gates. Give each leaf a
unary weight equal to its table entry when its indicator is one and to
$1$ otherwise. Exactly one leaf supplies the selected entry. Summing
the level costs gives $O(2^r)=O(D^2)$ per original station.
\end{proof}

\begin{proof}[Proof of Theorem~\ref{thm:shaving}]
Fix the finite field $k$ with $\omega_k<9/4$ and the exterior-bit
allowance $\eta>0$. Write $h=\lceil\log_2 n\rceil$. The ordinary-width
sort in Theorem~\ref{thm:sort} supplies an exponent $1<p<2$. Put
$\rho=p-1$. A total arity increase from $h^\beta$ to $h^{b_*}$ will
cost $O(h^{\rho b_*})$ per surviving incidence in the final lookup.
Choose $0<\delta\le1/10$ so that $1-\delta>\rho(1+\delta)$. Take
$A>0$ sufficiently small, and then $\beta>0$ sufficiently small, that
$b_*=\beta+(1+\delta)A$ satisfies
\begin{equation}\label{eq:shaving-margins}
 b_*<\tfrac12,\qquad \rho b_*<\tfrac12\eta,\qquad
 \Delta=(1-\delta)A-\rho b_*>0 .
\end{equation}
The positive gap $\Delta$ is the net saving on the original network
size. Choose a fixed integer $J$ large enough that $a=A/J$ satisfies
Lemma~\ref{lem:macro} uniformly for $\beta\le b\le1/2$ and
$2a<\eta/2$.

One extension $K/k$ suffices for these $J$ applications. Each width
scale $T$ lies between a constant times $h^\beta$ and a fixed power of
$h$. The local identities require a fixed constant extension and a
field of cardinality at least a fixed power of $T$. Choose $[K:k]$
large enough for the largest such power and divisible by the degrees
of the required constant extensions. Since $k$ is finite, we can do
this with $[K:k]=O(\log h)=O(\log T)$ at every scale.

Let $s_j$ count stations and their total arity after $j$ applications,
let $d_j$ be the cumulative number of copies, and put $e_j=d_j e$.
Initially $d_0=1$, $s_0=O(m)$, the arity is at most $3\le h^\beta$,
and the trivalent $k$-tables have a fixed finite dictionary. At step
$j<J$ use the size allowance
\[
 M_j=\max\{s_j,\lceil e_j h^{2a}\rceil\}.
\]
This allows the construction to continue even when the exterior bits
already account for most of the current size. We have $M_j\ge n$,
$e_j\le M_j h^{-2a}$ and maximum arity at most
$h^{\beta+(1+\delta)aj}$. If this application makes $r_j$ copies,
Lemma~\ref{lem:macro} preserves the short table labels, gives
$d_{j+1}=r_jd_j$ and yields
\[
 s_{j+1}\le r_j M_j q,\qquad q=h^{-(1-\delta)a}.
\]
Its construction in Section~\ref{sec:full} has
$\log_2 r_j=O((\log h)^2)$ uniformly in the input network. Inductively,
the original bound $s_0,e_0=O(n^{1.05})$, multiplied by a fixed
number of these copy factors and fixed powers of $h$, gives
$M_j\le n^{1.05+o(1)}\le n^{1.1}$ at every invocation.

In terms of size per exterior bit, $x_j=s_j/e_j$, the replication
factor cancels. For large $h$,
\[
 x_{j+1}\le q\max\{x_j,2h^{2a}\},\qquad
 x_J\le \max\{q^J x_0,2q h^{2a}\}.
\]
Indeed $q<1$, so every contribution arising from an exterior-bit
allowance is bounded by the most recent such allowance.

Let $u$ be the final maximum arity. Applying
Lemma~\ref{lem:batch-table} makes $2^u$ copies and multiplies the
size per copy by $O((1+u)^\rho)$. The resulting tensor is the
$d$th power of the original $k$-tensor, with
\[
 d=2^u d_J,\qquad
 \log_2 d\le h^{b_*}+O\bigl(J(\log h)^2\bigr)=o(h).
\]
Thus $d=n^{o(1)}$ uniformly at fixed $k,\eta$. Each of its $de$
exterior positions is retained through lookup and through
Lemma~\ref{lem:table-descent}. Descent costs $O([K:k]^2)=h^{o(1)}$.
Combining these bounds, the trivalent size over $k$ divided by $de$
is at most
\[
 O(m/e)\,h^{-\Delta+o(1)}
       +h^{2a+\rho b_*+o(1)} .
\]
Choose $\gamma=\Delta/2$. The strict inequalities
$\Delta>\gamma$ and $2a+\rho b_*<\eta$ absorb the constants and
lower-order terms for large $n$. Multiplication by $de$ gives
Equation~\eqref{eq:shave}.
\end{proof}

For the distinctness tensor of Theorem~\ref{thm:many}, the independent
exterior positions are the entries of separate lists. Choosing
$\eta<\sigma$, Proposition~\ref{prop:finitefield} iterates the
compression until the volume per exterior bit is $O(w^\eta)$, below
the required $w^\sigma$. This rules out $\omega_k<9/4$ over every
finite field. Lemma~\ref{lem:specialize} transfers the lower bound to
any field, where the matching upper bound gives Theorem~\ref{thm:square}.
